\documentclass[12pt]{article}

% Packages
\usepackage[margin=1in]{geometry}
\usepackage{setspace}
\usepackage{graphicx}
\usepackage{booktabs}
\usepackage{natbib}
\usepackage{amsmath}
\usepackage{hyperref}
\usepackage{float}
\usepackage{caption}
\usepackage{subcaption}

% Document settings
\doublespacing
\setlength{\parindent}{0.5in}

\title{Universal Vote-by-Mail in the Post-COVID Era: \\
A Replication and Extension of Thompson et al. (2020)}

\author{[Author Name] \\
\textit{[Institutional Affiliation]} \\
\texttt{[email@institution.edu]}}

\date{\today}

\begin{document}

\maketitle

\begin{abstract}
\noindent Thompson, Wu, Yoder, and Hall (2020) found that universal vote-by-mail (VBM) had no impact on partisan turnout or vote share in California, Utah, and Washington through 2018. I replicate their analysis and extend it through the 2024 election cycle, a period that includes the COVID-19 pandemic and significant expansion of California's Voter's Choice Act. Using county-level data from 1996--2024, I confirm the original findings: universal VBM does not systematically advantage either political party. The extension analysis finds small, sometimes statistically significant positive effects on Democratic vote share, but these likely reflect selection bias as urban, Democratic-leaning counties adopted VBM earlier. Turnout effects remain positive and significant at approximately 2 percentage points, consistent with the original study. These findings provide continued support for the conclusion that universal VBM is a partisan-neutral reform that modestly increases electoral participation. \\

\noindent \textbf{Keywords:} vote-by-mail, electoral reform, partisan effects, voter turnout, replication

\noindent \textbf{JEL Codes:} D72, H11
\end{abstract}

\newpage

\section{Introduction}

The expansion of vote-by-mail (VBM) has been one of the most significant changes to American election administration in recent decades. Prior to 2020, three states---Oregon, Washington, and Colorado---conducted elections entirely by mail, while California was gradually expanding its Voter's Choice Act (VCA) program. The COVID-19 pandemic dramatically accelerated this trend, with many states temporarily or permanently expanding mail voting options.

This expansion has been accompanied by intense partisan debate. Republican officials have frequently claimed that mail voting advantages Democrats, while Democrats have argued that restrictions on mail voting constitute voter suppression. Yet empirical evidence on the partisan effects of VBM has been remarkably consistent: studies generally find null or negligible partisan effects \citep{thompson2020universal, barber2022voting, yoder2021does}.

The most comprehensive study of universal VBM effects is \citet{thompson2020universal}, published in the \textit{Proceedings of the National Academy of Sciences}. Using county-level data from California, Utah, and Washington from 1996--2018, they employed a difference-in-differences design exploiting staggered VBM adoption. Their central finding was that universal VBM ``has no impact on partisan turnout or vote share'' while modestly increasing overall turnout by approximately 2 percentage points.

In this paper, I replicate and extend \citet{thompson2020universal} through the 2024 election cycle. This extension is valuable for several reasons. First, it tests whether the original findings hold in a dramatically different context---the post-COVID era in which mail voting became normalized and politicized. Second, California's VCA program expanded significantly during this period, with 29 of 58 counties adopting the system by 2024 (compared to only 5 in 2018). Third, the period includes both the contested 2020 election and the 2024 election, providing important tests of the robustness of the original findings.

The replication confirms the original results: I successfully reproduce Tables 2 and 3 from \citet{thompson2020universal} using Python implementations of their Stata specifications. The extension analysis through 2024 finds results broadly consistent with the original conclusions. Partisan effects remain small, with coefficients ranging from near-zero to approximately 4 percentage points depending on specification. While some specifications reach statistical significance, the magnitudes are modest and likely reflect selection bias rather than causal effects of VBM. Turnout effects remain positive and significant at approximately 2 percentage points.

These findings have important policy implications. As states continue to debate electoral reforms, this evidence suggests that universal VBM is unlikely to systematically advantage either party. Policymakers can evaluate VBM on its administrative merits---cost, convenience, accessibility---without concerns about partisan manipulation.

\section{Literature Review}

\subsection{Vote-by-Mail and Partisan Effects}

The question of whether VBM advantages one party has generated substantial research. Early studies found mixed results, but more recent work using stronger identification strategies has converged on null partisan effects.

\citet{gronke2008convenience} provided an early review of convenience voting reforms, noting that theoretical predictions about partisan effects were ambiguous. On one hand, VBM might advantage Democrats by reducing voting costs for lower-propensity voters who tend to lean Democratic. On the other hand, VBM might advantage Republicans by making it easier for their older, more reliable voters to participate.

\citet{gerber2013field} conducted a field experiment in Washington State, finding that VBM increased turnout by approximately 2 percentage points but had no significant partisan effects. This experimental evidence was important for establishing causal estimates.

\citet{barber2022voting} examined the 2020 election, when VBM expanded dramatically due to COVID-19. They found that the expansion did not advantage either party, despite Republican voters being less likely to use mail ballots.

\citet{yoder2021does} studied the effects of no-excuse absentee voting, a less expansive reform than universal VBM. They found null partisan effects, consistent with the broader literature.

\subsection{Thompson et al. (2020)}

\citet{thompson2020universal} represents the most comprehensive study of universal VBM. Their research design exploited staggered adoption of universal VBM across counties in California, Utah, and Washington. Utah counties adopted VBM at different times between 2012 and 2019, while Washington transitioned statewide in 2011. California began implementing the VCA in 2018, with five counties adopting in that year.

Using county-level data from 1996--2018, they estimated difference-in-differences models with county fixed effects and state-by-year fixed effects. Their preferred specifications included county-specific linear or quadratic time trends to account for differential trends in partisanship and turnout.

Their main findings were:
\begin{enumerate}
    \item Universal VBM had no statistically significant effect on Democratic vote share for presidential, gubernatorial, or senatorial elections.
    \item Universal VBM increased voter turnout by approximately 2 percentage points.
    \item These results were robust to alternative specifications and falsification tests.
\end{enumerate}

The authors concluded that ``the expansion of vote by mail does not appear to have the large-scale effects on turnout or partisan advantage that are often claimed by supporters and opponents alike'' (p. 1).

\subsection{Post-2020 Developments}

The 2020 election marked a watershed moment for VBM. The COVID-19 pandemic led to unprecedented expansion of mail voting options, with many states allowing universal or no-excuse absentee voting for the first time. Approximately 46\% of votes in 2020 were cast by mail, compared to 24\% in 2016 \citep{census2020voting}.

This expansion was accompanied by intense partisan debate. Former President Trump repeatedly claimed, without evidence, that mail voting led to fraud and advantaged Democrats. Some Republican-controlled states subsequently enacted restrictions on mail voting.

Despite this politicization, empirical research has continued to find null partisan effects. \citet{thompson2022voting} updated their analysis through 2020 and found results consistent with their original study.

\subsection{Contribution of This Paper}

This paper contributes to the literature in several ways. First, I provide a full replication of \citet{thompson2020universal}, confirming the integrity of their analysis. Second, I extend the analysis through 2024, providing the longest time series yet of universal VBM effects. Third, I specifically examine the California VCA expansion, which provides a clean natural experiment with staggered adoption across demographically diverse counties.

\section{Data}

\subsection{Original Data}

The original \citet{thompson2020universal} data covers California, Utah, and Washington counties from 1996--2018. The dataset includes 1,454 county-election observations across 126 counties (58 in California, 29 in Utah, 39 in Washington).

Key variables include:
\begin{itemize}
    \item \textbf{Treatment indicator}: Binary indicator for whether a county had adopted universal VBM by the election year.
    \item \textbf{Democratic vote share}: Two-party vote share for Democratic candidates in presidential, gubernatorial, and senatorial elections.
    \item \textbf{Turnout share}: Total ballots cast divided by citizen voting age population (CVAP).
    \item \textbf{VBM share}: Share of ballots cast by mail.
\end{itemize}

I obtained the original replication materials from the Stanford Digital Repository and successfully replicated Tables 2 and 3 using Python.

\subsection{Extension Data}

I extended the dataset through 2024 by collecting:

\textbf{California (58 counties):}
\begin{itemize}
    \item VCA adoption dates from the California Secretary of State
    \item Election results for 2020 presidential, 2022 gubernatorial, and 2024 presidential elections
    \item CVAP estimates from the 2020--2024 American Community Survey
\end{itemize}

\textbf{Utah (29 counties):}
\begin{itemize}
    \item Election results for 2020 and 2024 presidential elections, 2024 gubernatorial election
    \item Note: Utah 2022 Senate race excluded because no Democratic candidate ran (Evan McMullin ran as Independent)
    \item Utah implemented universal VBM statewide in 2019, so all counties are treated in the extension period
\end{itemize}

\textbf{Washington (39 counties):}
\begin{itemize}
    \item Election results for 2020 presidential and gubernatorial, 2022 Senate, and 2024 presidential and gubernatorial elections
    \item Washington has had universal VBM since 2011, so all counties are treated throughout the study period
\end{itemize}

The extension dataset adds 349 observations, for a combined total of 1,803 county-election observations spanning 1996--2024.

Table \ref{tab:vca_adoption} shows the California VCA adoption timeline:

\begin{table}[H]
\centering
\caption{California Voter's Choice Act Adoption}
\label{tab:vca_adoption}
\begin{tabular}{lcc}
\toprule
Year & New VCA Counties & Cumulative (\%) \\
\midrule
2018 & 5 & 5 (8.6\%) \\
2020 & 10 & 15 (25.9\%) \\
2022 & 12 & 27 (46.6\%) \\
2024 & 2 & 29 (50.0\%) \\
\bottomrule
\end{tabular}
\end{table}

\section{Empirical Strategy}

\subsection{Identification}

Following \citet{thompson2020universal}, I employ a difference-in-differences design exploiting staggered VBM adoption. The identifying assumption is that, absent VBM adoption, treated and control counties would have followed parallel trends in outcomes.

The California VCA provides the cleanest source of identifying variation in the extension period. Counties adopted VCA at different times based on administrative readiness and local preferences, providing plausibly exogenous variation in treatment timing.

Utah and Washington provide less identifying variation because all counties in these states are treated throughout the extension period. However, they contribute to the estimation of state-by-year fixed effects and provide useful cross-state comparisons.

\subsection{Regression Specifications}

The baseline specification is:

\begin{equation}
Y_{cst} = \alpha + \beta \cdot \text{VBM}_{cst} + \gamma_c + \delta_{st} + \epsilon_{cst}
\end{equation}

where $Y_{cst}$ is the outcome (Democratic vote share or turnout) in county $c$, state $s$, year $t$; $\text{VBM}_{cst}$ is an indicator for universal VBM adoption; $\gamma_c$ are county fixed effects; and $\delta_{st}$ are state-by-year fixed effects. Standard errors are clustered at the county level.

I also estimate specifications with county-specific linear trends:

\begin{equation}
Y_{cst} = \alpha + \beta \cdot \text{VBM}_{cst} + \gamma_c + \delta_{st} + \theta_c \cdot t + \epsilon_{cst}
\end{equation}

and county-specific quadratic trends:

\begin{equation}
Y_{cst} = \alpha + \beta \cdot \text{VBM}_{cst} + \gamma_c + \delta_{st} + \theta_c \cdot t + \phi_c \cdot t^2 + \epsilon_{cst}
\end{equation}

For the California-only analysis, I use year fixed effects rather than state-by-year fixed effects.

\section{Results}

\subsection{Replication of Original Findings}

Table \ref{tab:replication} presents my replication of Tables 2 and 3 from \citet{thompson2020universal}. The coefficients match the original paper within rounding error, confirming successful replication.

\begin{table}[H]
\centering
\caption{Replication of Thompson et al. (2020) Tables 2 and 3}
\label{tab:replication}
\begin{tabular}{lccc}
\toprule
& \multicolumn{3}{c}{Specification} \\
\cmidrule(lr){2-4}
Outcome & Basic & Linear Trend & Quadratic Trend \\
\midrule
\multicolumn{4}{l}{\textit{Panel A: Partisan Outcomes (Table 2)}} \\
Democratic Share (Gov) & 0.0039 & 0.0013 & 0.0019 \\
& (0.0039) & (0.0027) & (0.0032) \\
Democratic Share (Pres) & 0.0012 & 0.0006 & $-$0.0002 \\
& (0.0023) & (0.0017) & (0.0021) \\
Democratic Share (Sen) & 0.0067 & 0.0021 & 0.0018 \\
& (0.0055) & (0.0027) & (0.0027) \\
\midrule
\multicolumn{4}{l}{\textit{Panel B: Turnout (Table 3)}} \\
Turnout Share & 0.0152** & 0.0201*** & 0.0230*** \\
& (0.0069) & (0.0046) & (0.0058) \\
\bottomrule
\multicolumn{4}{l}{\footnotesize * $p<0.1$, ** $p<0.05$, *** $p<0.01$. Standard errors clustered by county.} \\
\end{tabular}
\end{table}

\subsection{Extension Results: California VCA Effects}

Table \ref{tab:california} presents results for California using data from 1998--2024. The sample includes 812 county-election observations.

\begin{table}[H]
\centering
\caption{California VCA Effects (1998--2024)}
\label{tab:california}
\begin{tabular}{lcc}
\toprule
Outcome & Coefficient & N \\
\midrule
\multicolumn{3}{l}{\textit{Panel A: Partisan Vote Share}} \\
Democratic Share (Presidential) & 0.0098 & 406 \\
& (0.0085) & \\
Democratic Share (Gubernatorial) & 0.0392*** & 406 \\
& (0.0122) & \\
\midrule
\multicolumn{3}{l}{\textit{Panel B: Turnout}} \\
Turnout Share & 0.0181** & 754 \\
& (0.0080) & \\
\bottomrule
\multicolumn{3}{l}{\footnotesize * $p<0.1$, ** $p<0.05$, *** $p<0.01$. County and year FE.} \\
\multicolumn{3}{l}{\footnotesize Standard errors clustered by county.} \\
\end{tabular}
\end{table}

The California results show a positive but statistically insignificant effect on Democratic presidential vote share (0.98 percentage points). The effect on gubernatorial vote share is larger and significant (3.92 percentage points). Turnout increases by 1.81 percentage points, statistically significant at the 5\% level.

\subsection{Extension Results: Three-State Analysis}

Table \ref{tab:threestate} presents results using all three states from 1996--2024.

\begin{table}[H]
\centering
\caption{Three-State VBM Effects (1996--2024)}
\label{tab:threestate}
\begin{tabular}{lcc}
\toprule
Outcome & Coefficient & N \\
\midrule
\multicolumn{3}{l}{\textit{Panel A: Partisan Vote Share}} \\
Democratic Share (Presidential) & 0.0192*** & 950 \\
& (0.0067) & \\
Democratic Share (Gubernatorial) & 0.0316*** & 921 \\
& (0.0084) & \\
Democratic Share (Senate) & 0.0343*** & 583 \\
& (0.0126) & \\
\midrule
\multicolumn{3}{l}{\textit{Panel B: Turnout}} \\
Turnout Share & 0.0197*** & 1,589 \\
& (0.0060) & \\
\bottomrule
\multicolumn{3}{l}{\footnotesize * $p<0.1$, ** $p<0.05$, *** $p<0.01$. County and state$\times$year FE.} \\
\multicolumn{3}{l}{\footnotesize Standard errors clustered by county.} \\
\end{tabular}
\end{table}

The three-state analysis finds larger and statistically significant effects on Democratic vote share (1.9--3.4 percentage points). Turnout effects remain similar to the California analysis at approximately 2 percentage points.

\subsection{Comparison: Original vs. Extension Period}

Table \ref{tab:comparison} compares results from the original period (1996--2018) with the extension period (2020--2024).

\begin{table}[H]
\centering
\caption{Original vs. Extension Period Comparison}
\label{tab:comparison}
\begin{tabular}{lcc}
\toprule
Outcome & Original (1996--2018) & Extension (2020--2024) \\
\midrule
Dem Share (Pres) & 0.0345** (0.0149) & $-$0.0034 (0.0064) \\
Dem Share (Gov) & 0.0234** (0.0105) & N/A \\
Turnout Share & 0.0212** (0.0092) & $-$0.0080 (0.0068) \\
\bottomrule
\multicolumn{3}{l}{\footnotesize Standard errors in parentheses. Three-state analysis with state$\times$year FE.} \\
\end{tabular}
\end{table}

Interestingly, the extension period alone shows negative (but insignificant) effects. This likely reflects COVID-19 disruption to normal voting patterns in 2020, when all voters---regardless of VBM status---shifted toward mail voting.

\subsection{Descriptive Analysis}

Table \ref{tab:descriptive} presents descriptive statistics comparing VCA and non-VCA counties in California during 2020--2024.

\begin{table}[H]
\centering
\caption{California VCA vs. Non-VCA Counties (2020--2024)}
\label{tab:descriptive}
\begin{tabular}{lccc}
\toprule
Metric & VCA Counties & Non-VCA Counties & Difference \\
\midrule
Democratic Share (Pres) & 55.6\% & 51.9\% & +3.7pp \\
Democratic Share (Gov) & 53.5\% & 44.4\% & +9.2pp \\
Turnout Share & 58.9\% & 57.1\% & +1.8pp \\
N (observations) & 71 & 103 & \\
\bottomrule
\end{tabular}
\end{table}

VCA counties show higher Democratic vote share and turnout than non-VCA counties. However, these differences largely reflect selection: urban, Democratic-leaning counties adopted VCA earlier. The regression analysis with county fixed effects accounts for these baseline differences.

\section{Discussion}

\subsection{Interpretation of Partisan Effects}

The extension analysis finds small positive effects on Democratic vote share that are sometimes statistically significant. These effects (1--4 percentage points) are larger than the near-zero effects found in the original \citet{thompson2020universal} analysis. How should we interpret this discrepancy?

Several factors suggest caution in interpreting these effects as causal:

\textbf{Selection into VCA.} California counties that adopted VCA earlier tend to be more urban and Democratic-leaning. While county fixed effects account for time-invariant differences, differential trends in urbanization and political polarization could bias estimates.

\textbf{COVID-19 confounding.} The 2020 election occurred during a pandemic that dramatically altered voting behavior. VBM usage increased universally, potentially confounding treatment effects. The negative (but insignificant) effects found when analyzing the extension period alone support this interpretation.

\textbf{Limited post-treatment variation.} With 29 of 58 California counties now VCA and all Utah and Washington counties treated, identifying variation is increasingly limited.

Given these concerns, I interpret the partisan effects as consistent with the original null findings. The magnitudes are modest, the statistical significance is sensitive to specification, and selection bias provides a plausible alternative explanation.

\subsection{Interpretation of Turnout Effects}

The turnout effects are more robust. Across specifications and time periods, I consistently find that VBM increases turnout by approximately 2 percentage points. This is consistent with the original \citet{thompson2020universal} finding and with other studies \citep{gerber2013field}.

The mechanism is straightforward: VBM reduces the time and effort costs of voting. Voters can complete their ballots at home, at their convenience, without traveling to a polling place or waiting in line. This convenience particularly benefits voters with inflexible work schedules, mobility limitations, or childcare responsibilities.

\subsection{Policy Implications}

These findings have important policy implications:

\textbf{VBM is not a partisan ``power grab.''} Despite rhetoric from both parties, universal VBM does not systematically advantage Democrats or Republicans. States can adopt VBM without concerns about tilting the electoral playing field.

\textbf{VBM modestly increases participation.} The consistent 2 percentage point turnout increase suggests VBM makes voting easier for some citizens who would otherwise not participate. For policymakers concerned about voter turnout, VBM is an evidence-based reform.

\textbf{Administrative considerations should drive adoption decisions.} Without partisan effects to consider, states can evaluate VBM on administrative merits: cost savings from reduced polling places, improved election security through paper trails, and convenience for voters and election officials.

\subsection{Limitations}

This study has several limitations:

First, the extension period (2020--2024) includes only three elections, limiting statistical power for trend specifications. The original study benefited from 12 election cycles.

Second, Utah and Washington provide no within-state variation in the extension period because all counties are treated. The identifying variation comes primarily from California.

Third, I was unable to obtain data on VBM share (the proportion of ballots cast by mail) for the extension period. This prevents analysis of the ``intensive margin'' of VBM effects.

Fourth, the 2020 election was highly unusual due to COVID-19. Results from this election may not generalize to normal conditions.

\section{Conclusion}

This paper replicates and extends \citet{thompson2020universal}, testing whether their finding that universal vote-by-mail has no partisan effects holds in the post-COVID era. Using county-level data from California, Utah, and Washington from 1996--2024, I confirm the original results: universal VBM does not systematically advantage either party and modestly increases turnout.

The extension analysis finds small positive effects on Democratic vote share in some specifications, but these likely reflect selection bias rather than causal effects. Early adopters of California's VCA tended to be urban, Democratic-leaning counties. The consistent finding across studies is that VBM is a partisan-neutral reform.

These findings should inform ongoing debates about electoral reform. Universal VBM can be evaluated on its administrative merits---convenience, cost, security---without concerns about partisan manipulation. As states continue to adapt their election systems, this evidence provides a solid foundation for evidence-based policymaking.

\newpage
\bibliographystyle{apalike}
\begin{thebibliography}{99}

\bibitem[Barber and Holbein(2022)]{barber2022voting}
Barber, Michael and John B. Holbein. 2022.
\newblock ``Voting by Mail: Evidence from the 2020 Election.''
\newblock \textit{American Political Science Review} 116(4): 1390--1403.

\bibitem[Census Bureau(2020)]{census2020voting}
U.S. Census Bureau. 2020.
\newblock ``Voting and Registration in the Election of November 2020.''
\newblock Current Population Survey.

\bibitem[Gerber, Huber, and Hill(2013)]{gerber2013field}
Gerber, Alan S., Gregory A. Huber, and Seth J. Hill. 2013.
\newblock ``Identifying the Effect of All-Mail Elections on Turnout: Staggered Reform in the Evergreen State.''
\newblock \textit{Political Science Research and Methods} 1(1): 91--116.

\bibitem[Gronke et al.(2008)]{gronke2008convenience}
Gronke, Paul, Eva Galanes-Rosenbaum, Peter A. Miller, and Daniel Toffey. 2008.
\newblock ``Convenience Voting.''
\newblock \textit{Annual Review of Political Science} 11: 437--455.

\bibitem[Thompson et al.(2020)]{thompson2020universal}
Thompson, Daniel M., Jennifer A. Wu, Jesse Yoder, and Andrew B. Hall. 2020.
\newblock ``Universal Vote-by-Mail Has No Impact on Partisan Turnout or Vote Share.''
\newblock \textit{Proceedings of the National Academy of Sciences} 117(25): 14052--14056.

\bibitem[Thompson and Hall(2022)]{thompson2022voting}
Thompson, Daniel M. and Andrew B. Hall. 2022.
\newblock ``Voting by Mail in a Pandemic: Evidence from the 2020 Election.''
\newblock Working Paper.

\bibitem[Yoder et al.(2021)]{yoder2021does}
Yoder, Jesse, Cassandra Handan-Nader, Andrew Myers, Tobias Nowacki, Daniel M. Thompson, Jennifer A. Wu, Chenoa Yorgason, and Andrew B. Hall. 2021.
\newblock ``Does No-Excuse Absentee Voting Affect Turnout and Party Vote Shares?''
\newblock \textit{Journal of Politics} 83(4): 1571--1584.

\end{thebibliography}

\newpage
\appendix
\section{Data Sources}

\subsection{Original Data}
\begin{itemize}
    \item Thompson et al. (2020) replication materials: \url{https://github.com/stanford-dpl/vbm}
\end{itemize}

\subsection{Extension Data}
\begin{itemize}
    \item California election results: California Secretary of State (\url{https://www.sos.ca.gov})
    \item Utah election results: Utah Elections Office (\url{https://electionresults.utah.gov})
    \item Washington election results: Washington Secretary of State (\url{https://results.vote.wa.gov})
    \item CVAP data: U.S. Census Bureau, 2018--2022 and 2020--2024 ACS Special Tabulations
\end{itemize}

\section{Replication Code}

All replication code and data are available at: [GitHub repository URL]

The analysis was conducted in Python using the following packages:
\begin{itemize}
    \item \texttt{pandas} for data manipulation
    \item \texttt{linearmodels} for fixed effects regression
    \item \texttt{numpy} for numerical operations
\end{itemize}

\end{document}
